\section{Discussion}\label{sec:discussion}

For top-of-line models with reasoning enabled, written guidance was enough: with reasoning \texttt{high}, the skill and
Agentomics-ML are indistinguishable on the test set for all four, and the skill costs a fraction as much. With
reasoning \texttt{none}, Agentomics-ML's fixed steps still help two of the four (Fable and Astra),
suggesting that its structure partly substitutes for the model planning its own work.

For three of the five mid-range models Agentomics-ML helps clearly: Haiku and DeepSeek score higher with it
at both reasoning settings, and Sonnet with reasoning \texttt{high}. The simplest explanation, that agents skip the skill's checks, does not hold: sixteen of the eighteen
skill runs ran the check script on every iteration and passed it (\cref{sec:limitations}). What differs is who
drives the procedure. Agentomics-ML gives each stage its own model turn and sends a script that fails back to the
step that wrote it; under the skill one model turn carries the whole task. The single-shot comparison
(\cref{tab:single}) points the same way: the state machine's advantage is largest on the first iteration, before the
agent has had a chance to correct its own course. What each method carries from one iteration to the next may explain
why the skill then catches up: the skill's agent keeps its whole conversation, while each state-machine step starts
afresh from a written summary of the earlier iterations (\cref{sec:conditions}). We did not test this directly.

Qwen 3.5 122B and GPT-4.1 are the exceptions, each in one of its two pairs. Qwen with reasoning \texttt{none} scored higher
with the skill; its state-machine run made no validation progress after its first iteration. GPT-4.1 scored
higher with the skill at both settings, completing three iterations where one of its state-machine runs failed the third; with
one run per setting the gap should not be over-read.

\paragraph{Reasoning and structure.} A state machine can be read as reasoning done in advance: the order of work and
the points at which to check it are decided before the model starts. If so, a model's own reasoning should matter more
under the skill, where that planning is left to it. The data do not settle it. With reasoning
\texttt{high} the skill's mean score rises by \SkillReasonGain{} (from \SkillNoneMean{} to \SkillHighMean), the state
machine's by \SmReasonGain{} (from \SmNoneMean{} to \SmHighMean); reasoning raised the score for \SkillHighWins{} of the 9 models
under the skill and \SmHighWins{} of the 9 under the state machine. The single-shot comparison shows the same pattern
more sharply: on a first attempt, before any feedback, the state machine's structure is worth the most.

The iteration analysis has a practical consequence: a two-iteration budget would have captured most of the
performance for about half the training time (\cref{tab:budget}).

\paragraph{What enforced structure buys.} A recurring question about harness-based agents is whether their structural
guarantees are worth the added complexity compared to a well-written skill or prompt. This experiment offers three observations
on that question. First, enforced structure matters for mid-range models: Sonnet, Haiku and DeepSeek all score
higher with Agentomics-ML, although the skill runs ran the same checks themselves. Second, no run could have seen the test labels while it was training, under either method. The skill makes the
test set a hand-over by the person: the agent finishes its iterations, reports, and only then is given the test
inputs, without labels, to predict on once. The state machine does the same in code: inputs and labels arrive after
the last iteration, for the harness's own final scoring step, when no model step remains; none of the tool calls those
runs made refers to the test upload. Both designs keep the test set out of reach by releasing it late;
the harness enforces the moment, and the skill leaves it to the person. Third, all reported scores were computed from the prediction files the runs wrote, not from anything the model
reported, so metric correctness is independently verified in both conditions. What the experiment does not address is whether
Agentomics-ML's step-wise context management helps or hurts; a direct comparison of full-context and step-wise
execution is a natural next step.

\paragraph{Cost.} The cost gap between the two conditions is large enough to matter in practice. The skill's eighteen
runs cost \SkillCost{} in total, Agentomics-ML's \SmCost{}. Much of the difference is structural: Agentomics-ML starts
a fresh model turn for each of its five model states per iteration, each carrying its own context, while the agent
works through the full task in a single turn. Within each condition, model tier is the dominant cost driver: top-of-line
models averaged \FlagSkillAvgCost{} per skill run and \FlagSmAvgCost{} per Agentomics-ML run, while mid-range models
averaged \SmallSkillAvgCost{} and \SmallSmAvgCost{} respectively. Fable 5.2 with reasoning \texttt{high} stands out at \FableSkillHighCost{} (skill) and \FableSmHighCost{}
(Agentomics-ML, over the two runs its record sums), driven by its large reasoning-token budget; DeepSeek V4.1 Flash, at the other extreme,
averaged under \$0.25 per run. For top-of-line models with reasoning enabled, where test performance is equivalent, the
skill is the more resource-efficient choice. For mid-range models, where Agentomics-ML scores higher, the extra cost
buys the enforced structure. Anyone choosing a configuration should weigh score and cost together.
